\documentclass[UTF8,11pt]{ctexart}
\usepackage[a4paper,top=2.6cm,bottom=2.6cm,left=2.8cm,right=2.8cm]{geometry}
\usepackage{array}
\usepackage{amsmath}
\usepackage{xcolor}
\setmainfont{Times New Roman}
\setCJKmainfont[BoldFont=SimHei]{SimSun}
\setmonofont{SimSun}
\setlength{\parskip}{3pt}
\begin{document}
\section*{\Large 第四部分：“盲区驭智”产品与技术解决方案}
\begin{quote}\small
编制说明：本章按申报书学术文体撰写。文中术语、公式、参数与结论均以项目现有资料（算法说明、仿真代码、开题报告、验证记录）为准；未完成或规划中的工作一律以“拟/计划/尚待验证”明确限定，不表述为已达成结论。文献引用编号见文末参考文献。
\end{quote}
---

\section*{一、研究定位}
\subsection*{（一）研究对象与系统定义}
大型货运车辆尤其是半挂车在低速右转过程中，由于前后轮轨迹半径不同，后轴内轮轨迹较前轴显著偏向转弯内侧，形成内轮差区域；对于牵引车—挂车铰接结构，挂车后轴无转向能力，其扫掠边界受车速、前轮转向角、铰接角、车辆尺寸及牵引销位置等多因素耦合影响，危险区域具有明显的动态演化特征[1-2][5]。现有后视镜、盲区镜及固定区域声光报警等装置，主要解决“盲区内是否存在目标”的感知问题，难以刻画“车辆未来短时内将扫掠何处”的时空冲突关系[3][6]。

针对上述问题，本项目将大型（铰接）车辆右转安全定义为\textbf{车辆右侧车身未来短时扫掠区域与弱势交通参与者运动状态之间的时空冲突预测问题}，构建一套以北斗高精度时空信息为基准、融合毫米波雷达与视觉感知、基于短时耦合运动学预测的实时预警系统。系统由车端嵌入式设备与仿真验证工具链两部分组成：车端设备完成车辆状态采集、扫掠危险区生成、目标风险评估与分级预警；工具链用于模型验证、参数整定与对照实验。

\subsection*{（二）典型应用场景与风险特征}
项目面向以下低速右转高频、车辆集中且管理主体明确的场景：

1. 城市交叉口：行人与非机动车高频混行，右转冲突风险集中[6]； 2. 港口与物流园区：半挂车、集装箱牵引车低速右转频繁，场地封闭，便于受控实验； 3. 工程运输：渣土车、搅拌车等右转事故多发，管理刚性强； 4. 干线货运与城市配送：半挂物流车队的中远期应用场景。

上述场景的共性风险特征为：车辆低速、铰接姿态持续变化、弱势交通参与者运动路径灵活，驾驶员难以仅凭视觉与经验判断右侧车身动态扫掠边界。

\subsection*{（三）系统功能定位与边界}
系统定位于\textbf{实时预警辅助}，其功能边界如下：

- 功能范围：实时采集车速 v、前轮转向角 α、铰接角 φ 及右侧目标信息；基于当前状态预测未来短时域右侧车身扫掠危险区；统一坐标后完成空间判定与风险评估，输出 0—3 级风险；分级声光报警与事件记录；增强版实现北斗轨迹与预警事件的时空关联； - 边界约束：本系统不涉及车辆纵向或横向控制，不生成控制指令；不做长时驾驶意图识别；不替代驾驶员；毫米波雷达有效探测范围有限（约 0—6 m），恶劣天气与极端遮挡下的性能尚待实车测试验证；TTC 复盘工具基于后验真值评估预警提前量，不参与在线判警。

---

\section*{二、系统总体技术架构}
系统采用“感知—时空基准—计算决策—输出”的分层体系，主循环周期 20 ms（50 Hz）。整体数据流为：车辆状态与外部目标信息经传感器层采集，与北斗时空基准共同输入嵌入式计算层；计算层完成耦合运动学预测、三层扫掠危险区生成、目标判内与风险分级；输出层执行声光报警与事件记录。

\begin{verbatim}
① 车辆状态感知层：车速（CAN）、前轮转向角（AS5600）、铰接角（角度电位计）
② 北斗高精度时空信息层（增强版）：位置、速度、航向
③ 毫米波雷达环境感知层：右侧目标距离、角度、速度
④ 视觉协处理感知层：行人、非机动车类别识别
⑤ 嵌入式计算与风险决策层：运动学预测 → 三层扫掠危险区 →
   目标判内 → 风险评估 → 报警状态机
⑥ 声光预警与数据记录层：分级报警、事件记录
\end{verbatim}
\subsection*{（一）车辆状态感知层}
\subsubsection*{1. 车辆速度采集}
系统通过 OBD/CAN 接口获取整车纵向车速 v。ESP32-S3 采用 TWAI 外设连接 CAN 收发器（SN65HVD230/TJA1050 级别），引脚 TX=GPIO5、RX=GPIO4，默认波特率 500 kbps。车速信号经滤波后作为运动学预测的纵向输入；CAN 报文解码器需在获取实车报文后完成，属第二阶段任务。

\subsubsection*{2. 前轮转向角采集}
前轮转向角 α 由 AS5600 磁性角度传感器测量，当前工程配置为模拟输出接入 ADC（GPIO34，12 位）。传感器零位、比例与方向标定参数已在固件中预留，需实测标定确定。

\subsubsection*{3. 牵引车—挂车铰接角采集}
铰接角 φ 由旋转角度电位计测量，模拟输出接入 ADC（GPIO35），零位标定同样预留待实测。

\subsection*{（二）北斗高精度时空信息层}
增强版接入北斗 RTK 模块，提供厘米级动态定位、高精度对地测速与双天线载波相位测向。该层为车体状态提供独立绝对真值，用于校正运动学积分产生的位置与航向漂移，并为“车体系—世界系”坐标变换提供统一基准。北斗经纬度按选定原点转换为本地 ENU 米制坐标，航向按“θ = 90° - 北向航向角”换算，以保证与算法坐标系一致。该层为增强版规划能力，基础版不安装，不影响在线判警主链路。

\subsection*{（三）毫米波雷达环境感知层}
毫米波雷达（HLK-LD2450，24 GHz）安装于车辆右侧，检测右侧约 0—6 m 范围内多目标的距离、角度与速度，经 UART（RX=GPIO16、TX=GPIO17，115200 baud）与主控通信。雷达帧解析器拟在第二阶段完成。

\subsection*{（四）视觉协处理感知层}
视觉协处理器（OpenMV H7 Plus）识别行人、非机动车等目标类别，默认 QVGA（320$\times$240），帧率低于 10 Hz 时降为 QQVGA；检测模型按轻量化约束设计（模型规模不大于 100 KB、单帧推理不超过 1 ms 级别）。视觉模块按固定二进制协议将目标类别、位置框与置信度传输给主控。该层完整实现属第三阶段任务。

\subsection*{（五）嵌入式计算与风险决策层}
主控 ESP32-S3（240 MHz）在每个 20 ms 周期内完成传感器读取、状态构造、三层扫掠危险区生成、目标判内（含当前车身占用区兜底）、全局风险评估与滞回状态机输出。按项目性能预算口径，核心算法单周期计算开销约 2 ms，多边形与目标缓冲区堆内存小于 5 KB；该指标为设计与仿真验证口径，最终实时性指标以正式实验记录为准。

\subsection*{（六）声光预警与数据记录层}
报警输出采用双色 LED（GPIO13）与无源蜂鸣器（GPIO12），分级输出为关闭、黄色预警、红色报警（蜂鸣器 250 ms 周期）、红色紧急（蜂鸣器连续）。报警状态机升级需条件连续满足 2 帧（40 ms），降级需条件连续失败 10 帧（200 ms），以抑制瞬时抖动误报。预警事件、车辆状态与目标信息同步记录，供离线复盘与后续数据服务使用。

---

\section*{三、半挂车辆短时耦合运动学预测}
\subsection*{（一）车辆运动状态建模}
以牵引车后轴中心为参考点建立车体坐标系（原点 B，x 轴向前，y 轴向左），车辆状态向量定义为：

\begin{equation}
s = [x_B, y_B, \theta, \varphi]^{\mathrm{T}}
\end{equation}
其中 x\_B、y\_B 为牵引车后轴中心位置（m），θ 为牵引车航向角（rad），φ 为牵引车—挂车铰接角（rad）。预测输入为车速 v（m/s）、前轮转向角 α（rad）与铰接角 φ（rad）。

车辆固定参数采用可配置结构，典型取值见表 4-1（参数随车型标定调整）。

\textbf{表 4-1 车辆固定参数（典型值）}

\begin{center}
\small
\begin{tabular}{p{3.2cm}p{7.4cm}p{2.6cm}}
\hline
\textbf{参数} & \textbf{含义} & \textbf{典型值} \\
\hline
l & 牵引车轴距（前轴 A → 后轴中心 B） & 4.0 m \\
\hline
l\_h & 后悬长（后轴 B → 鞍座牵引销 H） & 1.8 m \\
\hline
L & 挂车轴距（牵引销 H → 挂车后轴中心 T） & 13.5 m \\
\hline
W & 等效车宽 & 2.0 m \\
\hline
φ\_max & 铰接角硬限幅 & 15° \\
\hline
\end{tabular}
\end{center}
\begin{quote}\small
口径说明：典型值取自项目《预测与判警算法说明》；代码默认值与验证工况 CSV 数值可能因场景不同而调整（如示例工况曾采用 L=15 m、W=3 m、φ\_max=15°），实际以车辆参数标定为准。
\end{quote}
\subsection*{（二）牵引车—挂车耦合关系}
采用离散欧拉递推描述车辆姿态演化，一步递推关系为：

\begin{equation}
\omega_1 = \frac{v\sin\alpha}{l} 
\end{equation}
\begin{equation}
v_B = v\cos\alpha 
\end{equation}
\begin{equation}
\omega_2 = \frac{v_B\sin\varphi + l_h\omega_1\cos\varphi}{L} 
\end{equation}
\begin{equation}
v_T = v_B\cos\varphi - l_h\omega_1\sin\varphi 
\end{equation}
\begin{equation}
x_B \leftarrow x_B + v_B\cos\theta\,\mathrm{d}t,\quad y_B \leftarrow y_B + v_B\sin\theta\,\mathrm{d}t,\quad \theta \leftarrow \theta + \omega_1\,\mathrm{d}t,\quad \varphi \leftarrow \varphi + (\omega_1 - \omega_2)\,\mathrm{d}t,\quad |\varphi| \le \varphi_{\max} 
\end{equation}
其中 ω₁、ω₂ 分别为牵引车与挂车横摆率，v\_B 为牵引车后轴中心纵向速度，v\_T 为挂车纵向速度（备用）。式（3）保留 l\_h 耦合项，反映牵引销偏置对挂车横摆率的直接影响。

模型具有铰接角自回正性质：直线行驶（α=0，ω₁=0）时 dφ/dt = -v\_B·sinφ/L，铰接角自然衰减回零（v\_B=3 m/s 时时间常数约 4.5 s），与实车物理一致。该性质是判定模型几何正确性的依据，已在 MATLAB 与网页端以同方程实现并交叉验证。

\subsection*{（三）牵引销偏置与挂车内切特征}
半挂车右转危险区不能由单体车辆模型直接近似，机理有二：

其一，牵引销偏置耦合。牵引车转弯时，牵引销 H 点除随后轴平移外，还因后悬长度 l\_h 产生额外线速度，直接影响挂车横摆率；忽略该耦合项会使挂车轨迹推演产生偏差[2][4]。项目早期仿真（guacheweixianqu.m）即因缺失该耦合项导致 φ 不收敛，已修正并在代码中固化警示。

其二，挂车的内切与外甩特征。挂车后轴无转向能力，转弯时挂车尾部相对弯心向外甩出，而牵引销至挂车中部的右侧车身外缘向弯心内侧切入。实测反例（R3 目标）表明，真正形成“卷入—碾压”风险的主要是后者，即挂车中前段右侧车身向内侧扫掠形成的挤压区；尾部是否扫到目标取决于目标相对位置，不能一概而论。

\subsection*{（四）车辆未来短时姿态递推}
预测以当前状态为起点，按固定步长 dτ 递推至目标时间窗，窗口内前轮转向角默认采用一阶保持假设（α(τ)=α\_now）。代码保留线性外推能力（α(τ)=α₀+α̇·τ），但当前默认关闭：实测 PID 输出转向角含高频抖动，差分估计的 α̇ 噪声较大，长时外推易被放大导致危险区发散；后续启用需先低通滤波并限幅。当前 CFG 体系下预测步长统一为 0.02 s。

---

\section*{四、多时间尺度动态扫掠危险区}
\subsection*{（一）关键车辆几何点定义}
系统定义四个关键几何点表达整车轮廓（图略）：

\begin{equation}
A = B + l(\cos\theta,\ \sin\theta),\quad H = B + l_h(\cos\theta,\ \sin\theta),\quad \theta_t = \theta - \varphi,\quad T = H - L(\cos\theta_t,\ \sin\theta_t) 
\end{equation}
其中 A 为牵引车前轮中心，H 为鞍座牵引销，T 为挂车后轴中心，θ\_t 为挂车航向。

\subsection*{（二）车辆右侧边界计算}
挂车右侧外缘由 H—T 直棍向右法向偏移半车宽得到：

\begin{equation}
\boldsymbol{n}_r = (\sin\theta_t,\ -\cos\theta_t),\quad H_r = H + \frac{W}{2}\boldsymbol{n}_r,\quad T_r = T + \frac{W}{2}\boldsymbol{n}_r 
\end{equation}
预测中 A、B、H 点使用牵引车航向 θ 取右法线，T 点使用挂车航向 θ\_t 取右法线，避免以单一航向近似整车边界。

\subsection*{（三）0.3 秒紧急危险区}
0.3 s 窗口对应紧急危险区（PolyI），表达即将发生的直接碰撞风险。以 v=3 m/s 为例，纵向覆盖约 0.9 m，目标进入即按最高风险处理。

\subsection*{（四）1.0 秒报警危险区}
1.0 s 窗口对应报警危险区（PolyA），是在线 TTC 评估的基准区域。以 v=3 m/s 为例，纵向覆盖约 3 m，对应红色报警。

\subsection*{（五）2.0 秒预警危险区}
2.0 s 窗口对应预警危险区（PolyW），为最远预警窗口。以 v=3 m/s 为例，纵向覆盖约 6 m，对应黄色预警；盲区分析以 PolyW 首次命中与真实接触的时间差定义预警提前量。

构造方法上，系统逐预测步记录挂车右边缘两点，按“起始右缘—H\_r 弧线—结束右缘—T\_r 弧线—闭合”组成扫掠多边形。由于挂车尾部外甩，两条弧线可能交叉形成自相交多边形，采用偶奇规则拆分后仅保留弯心侧面积最大的危险区，剔除外甩三角。统计口径上，为避免长窗口合并为单一大多边形后自交或退化、导致“扩窗反而漏报”，盲区分析、MATLAB 校验与状态机仿真统一采用相邻预测步右边缘小扫掠四边形逐段判定（segment\_swept），保证扩窗单调性。

---

\section*{五、多传感器目标感知与信息融合}
\subsection*{（一）毫米波雷达目标位置与速度感知}
毫米波雷达提供右侧多目标的距离、角度与径向速度，输出为雷达安装坐标系下的表达。雷达数据经外参变换进入车体坐标系后方可参与判警；雷达帧解析拟在第二阶段完成。

\subsection*{（二）视觉目标类别识别}
视觉模块辅助回答“目标是什么”，识别行人、非机动车等类别（类别编码：0 未知、1 行人、2 自行车/电瓶车、3 汽车），输出目标框与置信度。视觉模型与识别链路为第三阶段规划任务。

\subsection*{（三）传感器外参标定}
多传感器协同的前提是统一坐标，系统规划三类标定：雷达安装位置与波束偏角标定（相对车体参考点）、视觉安装位姿外参标定、角度传感器零位与比例标定。标定方案已设计，实测标定随第二阶段台架测试实施。

\subsection*{（四）统一车体坐标系}
目标坐标经两级变换统一：

\begin{equation}
\boldsymbol{p}_B = \boldsymbol{R}(\psi_r)\boldsymbol{p}_R + \boldsymbol{t}_R,\qquad \boldsymbol{p}_W = \boldsymbol{R}(\theta)\boldsymbol{p}_B + [x_B,\ y_B]^{\mathrm{T}} 
\end{equation}
其中下标 R、B、W 分别表示雷达系、车体系与世界系，ψ\_r、t\_R 为雷达安装偏角与平移。在线判警将预测多边形预变换至车体系，使目标留在车体系内直接判内；世界系坐标用于轨迹记录、路口几何匹配与风险时空数据。

系统同时设置当前车身占用区（BodyNow）兜底判定：目标一旦落入当前时刻挂车车身矩形，即直接判定为最高风险。该机制修复了早期验证中“仅判未来扫掠、漏判已进入车身目标”的反例（R3 目标），是风险判定的必要组成部分。

\subsection*{（五）多源目标信息融合}
雷达目标（位置与速度）与视觉目标（类别与置信度）经关联后形成统一目标列表，采用短时恒速外推参与未来风险判定。第三阶段验收门控为：任一传感器漏检时，融合结果仍能恢复目标。融合实现为规划任务。

---

\section*{六、目标风险评估与分级预警}
\subsection*{（一）目标短时运动外推}
系统对每个目标按当前速度做短时恒速（CV）外推，用于未来时刻的空间判内与在线 TTC 评估（预测口径，见下文）。

\subsection*{（二）点在扫掠危险区内判定}
系统采用射线法实现点在多边形内判定。盲区分析中测试点常位于扫掠边界线上，故采用 0.05 m 几何容差，即目标在多边形内部或距边界不超过 5 cm 均视为命中。

\subsection*{（三）TTC 风险评估（双口径）}
项目中存在两种口径的 TTC，需严格区分：

- \textbf{复盘真值 TTC（后验真值）}：在已知整段行驶轨迹后，反推目标与车辆扫掠区域的真实碰撞时间（Time to Collision）。该值由 TTC 预警复盘软件基于完整真值轨迹计算，用于评估在线算法的预警提前量，不参与在线判警。项目算法说明明确“true TTC 只用于复盘评估主算法提前量”，在线端不能直接使用真值 TTC。 - \textbf{在线 TTC（预测口径）}：以目标短时恒速外推，求其最早进入报警危险区（PolyA）的时间，作为风险分级判据之一。按项目算法说明，该在线 TTC 链路与报警状态机滞回一并标注为第二阶段待实装；当前已实现并作为在线判警主判据的是“点在扫掠危险区多边形内判定 + 当前车身占用区兜底”。

\subsection*{（四）0—3 级风险分级}
\begin{center}
\small
\begin{tabular}{p{5.6cm}p{2.2cm}p{3.2cm}}
\hline
\textbf{判定条件} & \textbf{风险等级} & \textbf{报警输出} \\
\hline
落入预警区或在线 TTC<2.0 s & 1 & 黄色预警 \\
\hline
落入报警区或在线 TTC<1.0 s & 2 & 红色报警 \\
\hline
落入紧急区、在线 TTC<0.3 s 或命中当前车身占用区 & 3 & 紧急 \\
\hline
均不满足 & 0 & 关闭 \\
\hline
\end{tabular}
\end{center}
全局风险取所有目标风险的最大值，任一目标进入紧急状态即触发最高级报警。

\begin{quote}\small
口径说明：表中在线 TTC 阈值为规格设计并标注第二阶段待实装；当前在线判警主判据为点在多边形内判定与当前车身占用区兜底。复盘真值 TTC 仅用于后验评估预警提前量。
\end{quote}
\subsection*{（五）报警状态机与滞回机制}
为避免瞬时数据波动造成报警抖动，报警状态机采用滞回逻辑：升级需条件连续满足 2 帧（40 ms），降级需条件连续失败 10 帧（200 ms）。输出状态包括关闭、黄色预警、红色报警（蜂鸣 250 ms 周期）与红色紧急（蜂鸣连续）。该机制的嵌入式实装属第二阶段任务，当前已在仿真与工具链中完成设计验证。

---

\section*{七、北斗高精度定位的作用}
\subsection*{（一）车辆位置、速度与航向信息获取}
北斗高精度定位为系统提供三类独立时空信息：

- 位置：RTK 载波相位差分定位可达厘米级动态精度，用于替代或校正运动学积分位置（积分存在漂移）； - 速度：多普勒测速精度约 0.03—0.1 m/s，为对地速度、不受车轮打滑影响，可校正 CAN 车速； - 航向：双天线载波相位测向精度约 0.1°—0.5°（依基线长度）；单天线仅能测航迹角，车头航向需双天线或与惯性单元融合。

坐标层面，北斗经纬度按选定原点转换为本地 ENU 米制坐标，映射为算法统一的 x\_B、y\_B；航向按“θ = 90° - 北向航向角”换算。该层为增强版规划能力，模组型号待选型确定。

\subsection*{（二）预警事件时空关联}
北斗轨迹为预警事件补充时间与位置属性，使每次报警可关联至具体路口与轨迹片段，形成“何时、何地、何种风险等级”的事件记录，为车队管理提供事件复盘基础。

\subsection*{（三）车队风险数据沉淀}
汇聚多车预警事件与轨迹后，可形成车辆安全报告、风险路口统计、预警事件复盘等数据服务。该功能为规划方向，技术实现未承诺具体指标。

\subsection*{（四）路口风险黑点识别的拓展潜力}
基于预警事件时空聚类与路口几何关联，系统具备向路口风险黑点识别与道路治理延伸的潜力。该方向为拓展规划，未承诺具体算法实现或准确率指标。

---

\section*{八、系统软硬件实现}
\subsection*{（一）ESP32-S3 主控平台}
主控采用 ESP32-S3-N16R8（16 MB Flash、8 MB PSRAM，240 MHz），承载传感器读取与滤波、运动学预测、三层危险区生成、目标判内、风险状态机与报警输出。主循环按 20 ms（50 Hz）设计，核心算法预算约 2 ms/周期，堆内存占用小于 5 KB；以上为设计验证口径，最终以正式实验记录为准。

\subsection*{（二）OpenMV 视觉协处理模块}
OpenMV H7 Plus 经 UART2 与主控通信（115200 baud），默认 QVGA 分辨率，帧率低于 10 Hz 时降为 QQVGA；检测模型按轻量化约束（不大于 100 KB、单帧推理不超过 1 ms 级别）设计。

\subsection*{（三）毫米波雷达及角度传感器}
毫米波雷达（HLK-LD2450）负责右侧 0—6 m 目标检测（UART RX=16、TX=17）；AS5600 测量前轮转向角（GPIO34）；角度电位计测量铰接角（GPIO35）。传感器读数经 IIR 低通滤波后进入预测链路。

\subsection*{（四）CAN/OBD 车辆数据接口}
系统经 TWAI 外设与 CAN 收发器接入整车 CAN 总线（TX=GPIO5、RX=GPIO4，默认 500 kbps），通过 OBD-II 接口获取车速。CAN 报文解码需在实车报文基础上完成，属第二阶段任务。

\subsection*{（五）北斗 RTK 模块}
增强版配置支持双天线测向的双频 RTK 模组（具体型号待选型）、主/从 GNSS 天线与可选惯性单元，提供位置、速度、航向输入。基础版不安装该模块；增强版选型与 BOM 核算尚未定稿。

\subsection*{（六）系统通信与数据记录}
视觉模块与主控间采用固定二进制协议传输目标帧：

\begin{verbatim}
0xAA 0x55 N obj1 ... objN CRC8
每个 obj：cls(u8) x_mm(i16 LE) y_mm(i16 LE) w_mm(u16 LE) h_mm(u16 LE) conf(u8 0..100)
\end{verbatim}
目标坐标由视觉模块预做外参校准（车体系表达），帧率目标不低于 10 Hz。主控侧记录车辆状态、目标、风险等级与报警事件，供离线复盘；数据记录介质规划为存储卡与串口日志。

---

\section*{九、仿真与实验验证方案}
\subsection*{（一）工况生成}
PID 工况仿真导出器通过配置车辆参数、参考航向与 PID 增益，生成平滑可控的转弯工况，以 0.02 s 间隔输出车速、转向角、铰接角、航向及关键几何点序列，模拟 50 Hz 节拍下的传感器输入。PID 模型在该体系中定位为工况生成器与对照模型，不作为最终在线判警模型。

\subsection*{（二）运动学验证}
MATLAB 验证程序加载工况 CSV 后，从初始状态按同一运动学方程重算轨迹，与 CSV 真值对比，并生成三层扫掠危险区与误差报告。第一阶段验收门控为相同输入下重算轨迹与真值误差小于 1 mm（几何一致性口径）。在此基础上完成 alpha 预测假设的横向对比（一阶保持、线性外推、多假设并集、扇形包络、二阶外推）与正交路口出弯状态机仿真，用于“基线—改进”论证。

\subsection*{（三）盲区分析}
转弯全过程盲区分析软件导入场景与目标点后，按当前帧在线预测三层危险区，并计算边界点预警提前量：

\begin{verbatim}
提前量 = 真实接触时间 - 首次 PolyW 报警时间
提前量 < 反应阈值（默认 1.0 s）→ 判定为盲区点
\end{verbatim}
工具默认采用 0.05 m 判内容差、2.0 s 起始截断，统计口径为 segment\_swept。状态机阈值与分状态预测窗口由 CFG 策略文件统一控制（默认窗口 2/1/0.3 s，出弯阶段示例 3/1.5/0.5 s），并与 MATLAB 校验、阈值扫描保持同一口径。

\subsection*{（四）后验复盘}
TTC 预警复盘软件基于完整真值轨迹反推目标真实碰撞时间，与在线预测首次报警时间对比，量化预警提前量；该工具为后验评估，不作为在线判警依据。Python 回放工具提供交互式时间滑块查看车身、扫掠区与目标状态。规划中的人工驾驶仿真软件将支持人工驾驶、导出同格式 CSV、状态机相位日志与回放，用于评估策略在非平滑人工操作下的表现。

\subsection*{（五）半实物测试}
项目按“桌面测试—半实物台架—低速实车”顺序推进。半实物测试通过手动旋转转向角与铰接角、放置反射板，验证传感器标定、坐标统一与报警等级切换；第二阶段验收门控为台架输入下风险等级按预期升降。

\subsection*{（六）低速实车验证计划}
第四阶段计划开展低速（5—15 km/h）实车场测，记录预警提前量、误报率、漏报率、系统刷新周期、目标识别稳定性与多车型参数适配效率，并以 CSV 日志与报警复盘形成测试报告。在实车测试完成前，系统指标以仿真与台架结果为准，不以仿真结论替代实车结论。

---

\section*{十、系统形态与配置规划}
\subsection*{（一）基础车端预警版}
基础版面向后装市场，提供车辆状态采集、三层扫掠预测、雷达与视觉目标判警、分级声光报警与本地事件记录。按当前 BOM 估算，基础版硬件与材料成本约 1200 元（38 项 BOM 中的主控、传感、电源、接口、输出、结构与测试类）；该成本为初步经济性估算，具体以采购询价为准。

\subsection*{（二）北斗高精度增强版}
增强版在基础版上增加北斗 RTK 模组、双天线测向与可选惯性单元，提供位置、速度、航向基准与预警事件时空关联。增强部分 BOM 估算约 620 元；具体成本待器件选型与试点报价确定后重新测算。

\subsection*{（三）车队数据管理版}
面向规模化车队，提供北斗轨迹关联、风险事件统计、车辆安全报告、路口风险分析与远程升级运维服务。该功能为规划方向，相关技术指标待产品定义确定。

\subsection*{（四）OEM/算法集成方案}
面向挂车厂、商用车主机厂与汽车电子企业，提供算法 SDK、车型参数适配、嵌入式移植与联合开发服务，以物理模型可解释、轻量化易集成为技术基础。具体接口标准待前装合作需求确定。

---

\section*{十一、研究进展与后续计划}
\subsection*{（一）已完成内容}
1. 核心算法：耦合运动学模型、关键几何点与右边界计算、三层扫掠危险区（含自相交处理）、点判内、当前车身占用区兜底、0—3 级风险分级与滞回状态机设计； 2. 仿真验证工具链：PID 工况仿真导出器、MATLAB 运动学验证（误差门控）、alpha 预测假设对比、正交路口出弯状态机仿真、网页盲区分析软件、TTC 复盘软件、CFG 策略体系与配置生成器、MATLAB 校验与阈值扫描； 3. 对照实验结论：多假设并集方案平均扫掠面积由基线 4.56 m² 增至 6.88 m²，但对典型目标（R3）预警提前量无实质改善，故定位为保守并集对照；R3 反例经当前车身占用区兜底修复（真实车身覆盖时段 16.24—18.12 s 可正确判定）； 4. 文档与申报材料：开题报告、答辩材料、项目介绍、技术路线详解与 BOM 清单。

\subsection*{（二）正在验证内容}
1. 出弯扩窗策略：EXIT 阶段 3/1.5/0.5 s 与默认 2/1/0.3 s 的盲区点分布对比，三端口径已统一； 2. 人工驾驶数据：人工驾驶仿真软件处于规划与开发中，将产出人工驾驶 CSV 与相位日志，验证策略在非平滑操作下的表现； 3. 状态机参数整定：入弯、弯中、出弯、结束各阶段阈值与保持帧数在典型工况下的稳定性； 4. 传感器与坐标统一：雷达、视觉外参与角度传感器零位标定方案已设计，实测标定随台架测试实施。

\subsection*{（三）下一阶段研究任务}
\begin{center}
\small
\begin{tabular}{p{2.2cm}p{7.4cm}p{4.2cm}}
\hline
\textbf{阶段} & \textbf{主要任务} & \textbf{验收门控} \\
\hline
第二阶段 & ESP32 端预测、风险评估与报警状态机实装；雷达与 CAN 解码；角度传感器标定；半实物台架 & 台架输入下报警等级正确切换 \\
\hline
第三阶段 & OpenMV 视觉模型与通信链路；雷达—视觉融合关联 & 单传感器漏检时融合仍能恢复目标 \\
\hline
第四阶段 & 低速实车（5—15 km/h）测试与日志复盘；论文与知识产权申报 & 实车指标形成正式实验记录 \\
\hline
北斗增强 & RTK 模组选型、双天线测向、坐标转换与事件时空关联 & 增强版 BOM 与试点报价确定 \\
\hline
人工驾驶仿真 & 人工驾驶软件分阶段开发（物理模型、导出、回放、集成验证） & 导出 CSV 经现有工具链读取验证 \\
\hline
\end{tabular}
\end{center}
知识产权方面，项目核心技术、算法与软件均为团队独立研发，无知识产权争议；当前尚无已授权专利或软件著作权，计划在样机与工具链定型后，围绕短时扫掠危险区预测方法与多源时空融合预警装置推进软件著作权与实用新型专利布局。

---

\section*{参考文献}
[1] 李逸良, 邱信明. 车辆转弯时内轮差的运动学理论模型[J]. 力学与实践, 2017, 39(1): 94-99.

[2] 李英帅, 闫琦若, 甄云竹, 等. 铰链车右转内轮差区域范围确定方法[J]. 中国安全科学学报, 2022, 32(3): 103-108.

[3] 张卫华, 邱猛, 张祥松, 等. 半挂车转弯盲区实时监测预警与安全保障[J]. 中国安全科学学报, 2020, 30(8): 57-62.

[4] 杨云, 黄鹏飞, 周懿文, 等. 基于TruckSim半挂牵引车内轮差危险区模型研究[J]. 自动化与仪器仪表, 2019(10): 106-109.

[5] 褚正清, 刘家保, 宋星. 货车右转弯内轮差的建模研究[J]. 河北北方学院学报(自然科学版), 2015, 31(5): 49-52.

[6] 王清洲, 徐令选, 范鑫, 等. 货车右转视线盲区形成机制及交叉口安全性研究[J]. 中国安全科学学报, 2017, 27(5): 99-104.

[7] 陈思嘉. 汽车转弯时由内轮差引发的交通事故原因建模与分析[J]. 数学通报, 2013, 52(12): 55-57.

[8] 张正轩. 重卡驾驶员视野盲区监控系统研究[D]. 西安: 长安大学, 2017.

---

\section*{附：技术口径与事实来源说明}
1. 本文技术内容以项目仓库现有资料为准：AGENTS.md/CLAUDE.md、预测与判警算法说明、预测模型与对照模型构建思路、CFG 扩窗算法一致性说明、正交路口出弯状态机仿真说明、alpha 方案对比说明、开题报告、答辩 PPT、固件配置文件（config.h）、仿真代码（kinematics\_step.m、predict\_swept.m、vehicle\_params.m、load\_pid\_scenario.m）及论文引用汇总。 2. 已完成内容指已在仓库中实现或形成文档验证的内容；正在验证与规划内容分别以“正在验证”“拟/计划”明确标注。 3. 性能与精度类指标（20 ms 主循环、50 Hz、误差门控、实时预算等）为项目设计与仿真验证口径，最终以正式实验记录为准。 4. 成本与 BOM 数值为当前估算，采购与量产前需询价核价；北斗增强版成本待选型与 BOM 确定后重新测算。

\end{document}